% introduction.tex -- Introduction
\section{Introduction}

During the October 1987 market decline, portfolio insurance concentrated
selling across participating traders~\cite{shiller1988portfolio,grossman1988analysis}.
Forced sales by funds that had borrowed to invest spread losses across
quantitative equity strategies in August 2007~\cite{khandani2011what}.
Such episodes motivate a common modeling question: how can changes in
some investors' decision rules destabilize a market and contribute to
a crash? Agent-based
models address collective behavior through imitation, strategy adaptation,
and responses to a shared price~\cite{lebaron2006agent,hommes2006heterogeneous,samanidou2007agent,axtell2025agent}.
In the minority game, every trader reads the same record of past outcomes
and switches among strategies scored on that record, so traders interact
only through the market outcome~\cite{challet1997minority,challet2000modeling}.
In the Santa Fe artificial stock market, traders learn forecasting rules
from a shared price and dividend~\cite{lebaron1999time}; adding technical
rules to a trader's repertoire was a dominant strategy that lowered wealth
and raised volatility for all~\cite{joshi1999technical,joshi2002financial},
and how much technical trading emerged depended on details of the learning
algorithm~\cite{ehrentreich2006technical}.
Many of these models include technical traders, who decide from the recent
history of the price, alongside traders who follow fundamental value or the
opinions of others~\cite{brock1998heterogeneous,lux1999scaling,hommes2006heterogeneous}.
Herding models describe how individual decisions produce collective
activity~\cite{kirman1993ants,cont2000herd,lux1999scaling}, while lattice
models encode local interactions through stochastic imitation or spin-like
responses, sometimes alongside market-wide
signals~\cite{bornholdt2001ising,iori2002microsimulation,bartolozzi2004stochastic,qiu2007cellular,sornette2006importance}.
These interactions can generate large fluctuations and periods of sustained
volatility, two features of financial returns~\cite{cont2001empirical}.
These approaches explain fluctuations through several forms of interaction,
but their aggregate return statistics do not by themselves separate the
effects of a persistent local rule change from feedback through price.
Both the size of the changed group and the nature of the change matter. A temporary increase in selling
and a persistent change in how traders respond to their neighbors need
not produce the same collective outcome.

A cellular automaton represents investors as traders arranged on a square
grid, with each trader deciding whether to buy, hold, or sell from the
decisions of nearby traders. Each trader uses a lookup table to convert the
sum of its four neighbors' numbered decisions into its next decision, a
totalistic rule~\cite{wolfram1984universality,packard1985two,wolfram2002new}.
Wolfram's own market model placed buying and selling entities on a line,
each responding to its two neighbors under a fixed rule, and took the
running difference between buyers and sellers as the
price~\cite{wolfram2002new}. A recent model in which thirty-two agents
followed rules inspired by Rule~110, kept balance sheets, and set the price
from the counts of buy and sell positions produced series with the
statistics of real assets~\cite{valsan2022on}. Two-dimensional
cellular-automaton markets have instead used stochastic rules, typically a
probability that a trader imitates its neighbors, and varied that
probability or a market-wide factor to change market
behavior~\cite{wei2003cellular,bartolozzi2004stochastic,qiu2007cellular,fan2009effect,atman2012influence}.
The model studied here uses a deterministic table drawn from the complete
set of totalistic rules, so one table entry can change while the other
responses stay fixed. Asking a trader to sell once changes its current decision, whereas
changing the rule changes how it responds whenever the same neighbor sum
occurs again. The new rule remains in use even after the trader stops
selling. Starting two markets from the same history, then changing the rule
in only one, allows the effect of that change to be measured against
continued use of the original
rule~\cite{glasserman1992guidelines,law2015simulation}.
In this model, selected traders along the edge of the grid are
technical traders: they buy when the price falls below a band around its
recent average and sell when it rises above that band, and they are the only
traders whose decisions depend on the price. The question is how
many traders must receive the new rule before at least half of the interior
traders decide to sell, which we call a Sell majority, and whether this
instability can develop when the technical traders are kept in Hold, so that
trades still complete and the price still moves, but no decision responds to
it.

A market dominated by sellers and a market price crash are different
outcomes. Their connection depends partly on the resources needed to trade. A trader needs shares to sell or cash to
buy, and completing a trade requires someone willing and able to take the
other side. Dealers can complete orders that other traders do not match, but
their cash and share holdings limit how much they can buy and
sell~\cite{garman1976market,stoll1978supply,ho1981optimal,madhavan2000market,adrian2017dealer}.
Price-impact models describe how completed trades change the
price~\cite{kyle1985continuous,hasbrouck1991measuring,farmer2005predictive,bouchaud2009how}.
Agent-based markets have included limited trader resources, records of
outstanding buy and sell orders, borrowing limits, and dealers with limited
cash and
shares~\cite{raberto2001agent,chiarella2002simulation,thurner2012leverage,braunmunzinger2018agent}.
In the grid model, traders can keep deciding to sell even when they cannot
complete a sale, allowing instability in decisions to develop without a
price crash. Limits on completed trades also change the price that informs later decisions. Dealer capacity can therefore affect both selling decisions and prices. We coupled execution and price through a dealer whose finite cash and
share reserves retained a constant product, following the reserve-based
approach of constant-function markets~\cite{angeris2022constant}.
Each batch used a common execution price consistent with its net trade,
and the final reserve ratio set the next marginal quote. This construction
tied price movements to actual transfers and prevented an isolated round
trip that restored dealer inventory from extracting its cash. It also
made reserve depth determine both available funding and price sensitivity,
so a smaller dealer need not protect the market from a large decline.
The relevant comparison was how a given pattern of decisions traded
against different reserves, and how the resulting prices influenced
later decisions.

In this study, we used a three-state cellular automaton market to examine
how changes in investor decision rules produced market instability and
how trading resources connected those decisions to price declines.
We compared markets with a common history, then changed one response in
the lookup table of selected interior traders while preserving the
comparison market's rule. Keeping technical traders in Hold allowed the
same intervention to be studied without feedback from price, while still
calculating trades and financial outcomes. A one-time instruction to sell
and a screen of other single-response changes tested which features of
the intervention mattered. Comparisons across grid sizes and observation
periods assessed whether the fraction needed for a Sell majority depended
on the extent and duration of interaction. We also varied dealer reserves
and changed trader cash and shares separately to identify their effects
on completed orders, majority formation, and drawdown. The paired design
therefore addressed two linked questions: whether a persistent change in
local decisions could spread through neighboring traders, and how the
execution mechanism determined the accompanying price decline.
